\documentclass[a4paper,12pt,twoside]{memoir}

% Castellano
\usepackage[spanish,es-tabla]{babel}
\selectlanguage{spanish}
\usepackage[utf8]{inputenc}
\usepackage[T1]{fontenc}
\usepackage{lmodern} % Scalable font
\usepackage{microtype}
\usepackage{placeins}

\RequirePackage{booktabs}
\RequirePackage[table]{xcolor}
\RequirePackage{xtab}
\RequirePackage{multirow}

% Links
\PassOptionsToPackage{hyphens}{url}\usepackage[colorlinks]{hyperref}
\hypersetup{
	allcolors = {red}
}

% Ecuaciones
\usepackage{amsmath}

% Rutas de fichero / paquete
\newcommand{\ruta}[1]{{\sffamily #1}}

% Párrafos
\nonzeroparskip

% Huérfanas y viudas
\widowpenalty100000
\clubpenalty100000

% Imágenes

% Comando para insertar una imagen en un lugar concreto.
% Los parámetros son:
% 1 --> Ruta absoluta/relativa de la figura
% 2 --> Texto a pie de figura
% 3 --> Tamaño en tanto por uno relativo al ancho de página
\usepackage{graphicx}
\newcommand{\imagen}[3]{
	\begin{figure}[!h]
		\centering
		\includegraphics[width=#3\textwidth]{#1}
		\caption{#2}\label{fig:#1}
	\end{figure}
	\FloatBarrier
}

% Comando para insertar una imagen sin posición.
% Los parámetros son:
% 1 --> Ruta absoluta/relativa de la figura
% 2 --> Texto a pie de figura
% 3 --> Tamaño en tanto por uno relativo al ancho de página
\newcommand{\imagenflotante}[3]{
	\begin{figure}
		\centering
		\includegraphics[width=#3\textwidth]{#1}
		\caption{#2}\label{fig:#1}
	\end{figure}
}

% El comando \figura nos permite insertar figuras comodamente, y utilizando
% siempre el mismo formato. Los parametros son:
% 1 --> Porcentaje del ancho de página que ocupará la figura (de 0 a 1)
% 2 --> Fichero de la imagen
% 3 --> Texto a pie de imagen
% 4 --> Etiqueta (label) para referencias
% 5 --> Opciones que queramos pasarle al \includegraphics
% 6 --> Opciones de posicionamiento a pasarle a \begin{figure}
\newcommand{\figuraConPosicion}[6]{%
  \setlength{\anchoFloat}{#1\textwidth}%
  \addtolength{\anchoFloat}{-4\fboxsep}%
  \setlength{\anchoFigura}{\anchoFloat}%
  \begin{figure}[#6]
    \begin{center}%
      \Ovalbox{%
        \begin{minipage}{\anchoFloat}%
          \begin{center}%
            \includegraphics[width=\anchoFigura,#5]{#2}%
            \caption{#3}%
            \label{#4}%
          \end{center}%
        \end{minipage}
      }%
    \end{center}%
  \end{figure}%
}

%
% Comando para incluir imágenes en formato apaisado (sin marco).
\newcommand{\figuraApaisadaSinMarco}[5]{%
  \begin{figure}%
    \begin{center}%
    \includegraphics[angle=90,height=#1\textheight,#5]{#2}%
    \caption{#3}%
    \label{#4}%
    \end{center}%
  \end{figure}%
}
% Para las tablas
\newcommand{\otoprule}{\midrule [\heavyrulewidth]}
%
% Nuevo comando para tablas pequeñas (menos de una página).
\newcommand{\tablaSmall}[5]{%
 \begin{table}
  \begin{center}
   \rowcolors {2}{gray!35}{}
   \begin{tabular}{#2}
    \toprule
    #4
    \otoprule
    #5
    \bottomrule
   \end{tabular}
   \caption{#1}
   \label{tabla:#3}
  \end{center}
 \end{table}
}

%
% Nuevo comando para tablas pequeñas (menos de una página).
\newcommand{\tablaSmallSinColores}[5]{%
 \begin{table}[H]
  \begin{center}
   \begin{tabular}{#2}
    \toprule
    #4
    \otoprule
    #5
    \bottomrule
   \end{tabular}
   \caption{#1}
   \label{tabla:#3}
  \end{center}
 \end{table}
}

\newcommand{\tablaApaisadaSmall}[5]{%
\begin{landscape}
  \begin{table}
   \begin{center}
    \rowcolors {2}{gray!35}{}
    \begin{tabular}{#2}
     \toprule
     #4
     \otoprule
     #5
     \bottomrule
    \end{tabular}
    \caption{#1}
    \label{tabla:#3}
   \end{center}
  \end{table}
\end{landscape}
}

%
% Nuevo comando para tablas grandes con cabecera y filas alternas coloreadas en gris.
\newcommand{\tabla}[6]{%
  \begin{center}
    \tablefirsthead{
      \toprule
      #5
      \otoprule
    }
    \tablehead{
      \multicolumn{#3}{l}{\small\sl continúa desde la página anterior}\\
      \toprule
      #5
      \otoprule
    }
    \tabletail{
      \hline
      \multicolumn{#3}{r}{\small\sl continúa en la página siguiente}\\
    }
    \tablelasttail{
      \hline
    }
    \bottomcaption{#1}
    \rowcolors {2}{gray!35}{}
    \begin{xtabular}{#2}
      #6
      \bottomrule
    \end{xtabular}
    \label{tabla:#4}
  \end{center}
}

%
% Nuevo comando para tablas grandes con cabecera.
\newcommand{\tablaSinColores}[6]{%
  \begin{center}
    \tablefirsthead{
      \toprule
      #5
      \otoprule
    }
    \tablehead{
      \multicolumn{#3}{l}{\small\sl continúa desde la página anterior}\\
      \toprule
      #5
      \otoprule
    }
    \tabletail{
      \hline
      \multicolumn{#3}{r}{\small\sl continúa en la página siguiente}\\
    }
    \tablelasttail{
      \hline
    }
    \bottomcaption{#1}
    \begin{xtabular}{#2}
      #6
      \bottomrule
    \end{xtabular}
    \label{tabla:#4}
  \end{center}
}

%
% Nuevo comando para tablas grandes sin cabecera.
\newcommand{\tablaSinCabecera}[5]{%
  \begin{center}
    \tablefirsthead{
      \toprule
    }
    \tablehead{
      \multicolumn{#3}{l}{\small\sl continúa desde la página anterior}\\
      \hline
    }
    \tabletail{
      \hline
      \multicolumn{#3}{r}{\small\sl continúa en la página siguiente}\\
    }
    \tablelasttail{
      \hline
    }
    \bottomcaption{#1}
  \begin{xtabular}{#2}
    #5
   \bottomrule
  \end{xtabular}
  \label{tabla:#4}
  \end{center}
}



\definecolor{cgoLight}{HTML}{EEEEEE}
\definecolor{cgoExtralight}{HTML}{FFFFFF}

%
% Nuevo comando para tablas grandes sin cabecera.
\newcommand{\tablaSinCabeceraConBandas}[5]{%
  \begin{center}
    \tablefirsthead{
      \toprule
    }
    \tablehead{
      \multicolumn{#3}{l}{\small\sl continúa desde la página anterior}\\
      \hline
    }
    \tabletail{
      \hline
      \multicolumn{#3}{r}{\small\sl continúa en la página siguiente}\\
    }
    \tablelasttail{
      \hline
    }
    \bottomcaption{#1}
    \rowcolors[]{1}{cgoExtralight}{cgoLight}

  \begin{xtabular}{#2}
    #5
   \bottomrule
  \end{xtabular}
  \label{tabla:#4}
  \end{center}
}



\graphicspath{ {./img/} }

% Capítulos
\chapterstyle{bianchi}
\newcommand{\capitulo}[2]{
	\setcounter{chapter}{#1}
	\setcounter{section}{0}
	\setcounter{figure}{0}
	\setcounter{table}{0}
	\chapter*{\thechapter.\enskip #2}
	\addcontentsline{toc}{chapter}{\thechapter.\enskip #2}
	\markboth{#2}{#2}
}

% Apéndices
\renewcommand{\appendixname}{Apéndice}
\renewcommand*\cftappendixname{\appendixname}

\newcommand{\apendice}[1]{
	%\renewcommand{\thechapter}{A}
	\chapter{#1}
}

\renewcommand*\cftappendixname{\appendixname\ }

% Formato de portada
\makeatletter
\usepackage{xcolor}
\newcommand{\tutor}[1]{\def\@tutor{#1}}
\newcommand{\course}[1]{\def\@course{#1}}
\definecolor{cpardoBox}{HTML}{E6E6FF}
\def\maketitle{
  \null
  \thispagestyle{empty}
  % Cabecera ----------------
\noindent\includegraphics[width=\textwidth]{cabecera}\vspace{1cm}%
  \vfill
  % Título proyecto y escudo informática ----------------
  \colorbox{cpardoBox}{%
    \begin{minipage}{.8\textwidth}
      \vspace{.5cm}\Large
      \begin{center}
      \textbf{TFG del Grado en Ingeniería Informática}\vspace{.6cm}\\
      \textbf{\LARGE\@title{}}
      \end{center}
      \vspace{.2cm}
    \end{minipage}

  }%
  \hfill\begin{minipage}{.20\textwidth}
    \includegraphics[width=\textwidth]{escudoInfor}
  \end{minipage}
  \vfill
  % Datos de alumno, curso y tutores ------------------
  \begin{center}%
  {%
    \noindent\LARGE
    Presentado por \@author{}\\ 
    en Universidad de Burgos --- \@date{}\\
    Tutor: Dr. \@tutor{}\\
  }%
  \end{center}%
  \null
  \cleardoublepage
  }
\makeatother

\newcommand{\nombre}{Adrián Zamora Sánchez} %%% cambio de comando

% Datos de portada
\title{Diseño e implementación de un lenguaje de programación y su compilador}
\author{\nombre}
\tutor{César Ignacio García Osorio}
\date{\today}

\begin{document}

\maketitle


\newpage\null\thispagestyle{empty}\newpage


%%%%%%%%%%%%%%%%%%%%%%%%%%%%%%%%%%%%%%%%%%%%%%%%%%%%%%%%%%%%%%%%%%%%%%%%%%%%%%%%%%%%%%%%
\thispagestyle{empty}


\noindent\includegraphics[width=\textwidth]{cabecera}\vspace{1cm}

\noindent D. nombre tutor, profesor del departamento de nombre departamento, área de Lenguajes y Sistemas Informáticos.

\noindent Expone:

\noindent Que el alumno D. \nombre, con DNI 48838775T, ha realizado el Trabajo final de Grado en Ingeniería Informática titulado título de TFG. 

\noindent Y que dicho trabajo ha sido realizado por el alumno bajo la dirección del que suscribe, en virtud de lo cual se autoriza su presentación y defensa.

\begin{center} %\large
En Burgos, {\large \today}
\end{center}

\vfill\vfill\vfill

% Author and supervisor
\begin{minipage}{0.45\textwidth}
\begin{flushleft} %\large
% para casos con solo un tutor comentar lo anterior
% y descomentar lo siguiente
Vº. Bº. del Tutor:\\[2cm]
César Ignacio García Osorio
\end{flushleft}
\end{minipage}
\hfill
\begin{minipage}{0.45\textwidth}
\begin{flushleft} %\large
%Vº. Bº. del co-tutor:\\[2cm]
%D. nombre co-tutor
\end{flushleft}
\end{minipage}
\hfill

\vfill

\newpage\null\thispagestyle{empty}\newpage




\frontmatter

% Abstract en castellano
\renewcommand*\abstractname{Resumen}
\begin{abstract}
En este proyecto se aborda la creación de un lenguaje de programación, herramienta fundamental dentro del campo de la ingeniería informática. Para que este lenguaje sea utilizable, se implemnetará un compilador propio, capaz de llevar este lenguaje a ejecución. Este proyecto se limitará a algunas funcionalidades clave, puesto que querer abordar cuestiones como OOP podrían aumentar la complejidad del proyecto demasiado.

Para la realización del proyecto se tratan de utilizar las herramientas más modernas y extendidas en la creación de lenguajes y compiladores: ANTLR, LLVM, CMake y C++ como lenguaje de programación. 
\end{abstract}

\renewcommand*\abstractname{Descriptores}
\begin{abstract}
Compilador, lenguaje de programación, ANTLR, LLVM
\end{abstract}

\clearpage

% Abstract en inglés
\renewcommand*\abstractname{Abstract}
\begin{abstract}
This project focuses on the creation of a programming language, a fundamental tool within the field of computer engineering. To make this language usable, a custom compiler will be implemented, capable of executing programs written in it. The scope of the project will be limited to a set of key functionalities, as tackling more advanced features such as OOP could increase the complexity beyond what is feasible.

For the development of the project, the most modern and widely adopted tools in language and compiler design will be used: ANTLR, LLVM, CMake, and C++ as the programming language.
\end{abstract}

\renewcommand*\abstractname{Keywords}
\begin{abstract}
Compiler, programmin language, ANTLR, LLVM
\end{abstract}

\clearpage

% Indices
\tableofcontents

\clearpage

\listoffigures

\clearpage

\listoftables
\clearpage

\mainmatter
\capitulo{1}{Introducción}
En la informática moderna estamos acostumbrados a disponer de numerosos lenguajes de programación, los cuales empleamos sin detenernos a analizar más allá de sus potenciales usos y librerías disponibles. Sin embargo, estas complejas herramientas cuentan con un gran trabajo detrás, el cual permite a los programadores utilizar estas herramientas sin preocuparse de cómo funcionan realmente.

La creación de los compiladores es una ciencia que ha evolucionado mucho a lo largo de los años, no solo por el surgimiento de nuevos lenguajes y sus compiladores asociados, muchos lenguajes que llevan con nosotros desde los 70s han experimentado cambios significativos en sus compiladores, además de versiones alternativas, cada una enfocanda en unos objetivos diferentes.

Ejemplos de proyectos que comenzaron siendo pequeños lenguajes de programación o lenguajes de dominio concreto que han crecido hasta ser herramientas ampliamente utilizadas son: 
\begin{itemize}
    \item \textbf{HTML + CSS}: Nacieron como lenguajes muy específicos que hoy día son estándares globales, tienen un ecosistema inmenso, disponen de importantes frameworks e incluso han inspirado a lenguajes y herramientas fuera del entorno web.
    \item \textbf{MATLAB}: Diseñado en los 80s para ser un lenguaje específico de algebra lineal, que hoy día se utiliza ampliamente en ingeniería, robótica, procesamiento de señales, visión por computadora, etc.  
    \item \textbf{SQL}: Comenzó en los 70s como un lenguaje que IBM desarrollaría para hacer consultas en sus bases de datos relacionales, sin embargo, fue ampliamente adoptado como un estándar que lo ha llevado a estar presente en prácticamente todos los motores de bases de datos relaciones modernos. 
\end{itemize}

\subsection{Memoria}
Esta memoria aborda los siguientes apartados:
\begin{itemize}
    \item \textbf{Objetivos}:
    \item \textbf{Conceptos teóricos}:
    \item \textbf{Técnicas y herramientas}:
    \item \textbf{Aspectos relevantes del desarrollo del proyecto}:
    \item \textbf{Trabajos relacionados}:
    \item \textbf{Conclusiones y líneas de trabajo futuras}:
\end{itemize}

\subsection{Anexos}
Adicionalmente, se adjuntan los siguientes anexos:
\begin{itemize}
    \item \textbf{Anexo A. Plan de proyecto}:
    \item \textbf{Anexo B. Requisitos}:
    \item \textbf{Anexo C. Diseño}:
    \item \textbf{Anexo D. Manual de programador}:
    \item \textbf{Anexo E. Manual de usuario}:
    \item \textbf{Anexo F. Sostenibilización curricular}:
\end{itemize}
\capitulo{2}{Objetivos del proyecto}

\begin{itemize}
    \item \textbf{Diseñar un lenguaje de programación}: 
    Definir la sintaxis y semántica de un nuevo lenguaje, estableciendo sus estructuras básicas, tipos de datos, operadores y reglas gramaticales.

    \item \textbf{Desarrollar un compilador capaz de reconocer el lenguaje}: 
    Implementar un compilador que analice y procese el código fuente del lenguaje, abarcando las fases de análisis léxico, sintáctico y semántico.

    \item \textbf{Generación y gestión de errores}: 
    Incorporar un sistema de detección y reporte de errores léxicos, sintácticos y semánticos, proporcionando mensajes claros y detallados que faciliten la depuración del código.

    \item \textbf{Herramienta de análisis y depuración visual}: 
    Desarrollar una utilidad visual que permita representar y explorar estructuras internas del compilador, como el Árbol de Sintaxis Abstracta (AST), facilitando la comprensión y depuración del lenguaje.

    \item \textbf{Compilación a partir de una representación intermedia (IR)}: 
    Implementar la traducción del lenguaje diseñado a una representación intermedia basada en LLVM IR, permitiendo la optimización y posterior generación de código máquina para distintas arquitecturas.

    \item \textbf{Implementación de utilidades internas del lenguaje}: 
    Desarrollar un conjunto de funciones o bibliotecas básicas que proporcionen utilidades estándar dentro del entorno del lenguaje diseñado.
\end{itemize}

\subsection{Objetivos técnicos}

El desarrollo del proyecto se apoya en el uso de herramientas y tecnologías ampliamente reconocidas en el ámbito de los compiladores modernos:

\begin{itemize}
    \item \textbf{ANTLR}: Se utilizará ANTLR como generador de analizadores léxicos y sintácticos, dada su modernidad, robustez y amplia adopción en el desarrollo de lenguajes formales.

    \item \textbf{LLVM}: Se empleará LLVM como infraestructura de compilación, debido a su capacidad para generar una representación intermedia (IR) portable y optimizable, independiente de la arquitectura de hardware.

    \item \textbf{C++}: El proyecto se desarrollará en C++, por su excelente integración tanto con ANTLR (a través del ANTLR Runtime para C++) como con LLVM, además de sus ventajas en programación orientada a objetos y gestión eficiente de memoria.

    \item \textbf{LaTex forest}: Esta herramienta ampliamente utilizada en el ámbio académico se utiliza para visualizaciones y diagramas. 
\end{itemize}
\capitulo{3}{Conceptos teóricos}

\subsection{Sobre lenguajes}
\begin{description}
    \item [Clasificación de lenguajes]: Los lenguajes de programación suelen clasificarse según su nivel de abstracción (alto o bajo nivel), paradigma (imperativo, declarativo, OOP, funcional), propósito del lenguaje (propósito general, lenguaje específico del dominio) y por forma de ejecución (compilados o interpretados).  
    \item [Tipado]: El tipado de un lenguaje se refiere a la forma en la que se gestionan los tipos de los datos, en los lenguajes con tipado estático, los tipos deben ser definidos en las declaraciones y se hacen comprobaciones estrictas durante las asignaciones, por otro lado, los lenguajes de tipado dinámico permiten declaraciones sin especificar el tipo, en estos lenguajes normalmente se hacen transformaciones de datos en tiempo de ejecución. 
    \item [Lenguaje máquina]: El lenguaje máquina es aquel que puede ser ejecutado por un ordenador, este tipo de lenguaje suele ser generado por un compilador, puesto que sería extremadamente ineficiente de escribir por un ser humano.
    \item [Gramaticas formales]: Las gramáticas formales son estructuras de reglas con las cuales se pueden generar las cadenas de caracteres permitidas por un lenguaje formal, esta gramática no tiene en cuenta el significado de las fórmulas bien formadas, únicamente la admisión de su forma.
\end{description}

\subsection{Sobre compiladores}
\begin{description}
    \item [Compilador por fases]: Es una técnica de construcción de compiladores que destaca por su alta modularidad, cada fase está bien definida y tiene una única responsabilidad, se caracterizan por ser fácilmente modificables y ampliables. 
    \item [Análisis léxico]: Es una fase del proceso de compilación en la cual se analizan las secuencias de caracteres de un texto y se separan en tokens o lexemas. A la herramienta que realiza este proceso se le llama scaner, lexer o tokenizer
    \item [Análsis sintáctico]: Esta fase del proceso de compilación trata de agrupar los tokens obtenidos durante el análisis léxico, para ello se vale de reglas que generan otras estructuras de datos (como AST) desde la cadena de tokens.
    \item [AST]: el Abstract Syntax Tree (árbol de sintaxis abstracta), es una estructura de datos que se utiliza comúnmente en compiladores para representar producciones sintácticas, donde los tokens forman una estructura gerárquica de árbol que representa el programa y ayuda a su posterior interpretación.
    \item [Tabla de símbolos]: Es una estructura de datos utilizada en compiladores para asociar cada símbolo de un programa con su ubicación, alcance y tipo de dato. Dentro de un compilador cumple un papel fundamental entre el front end y el back end del compilador.
    \item [Análisis semántico]: Es una fase en la cual el compilador comprueba la corrección de las producciones válidas formadas en el análisis sintáctico, algunas correcciones podrían ser la verificación de tipos y corrección en asignaciones y expresiones. Normalmente en esta fase se separa la información necesaria para rellenar la tabla de símbolos con los datos de los identificadores presentes en el código.
    \item [Generación de código intermedio]: la creación de una representación intermedia o IR por sus siglas en inglés, hace referencia a un código que queda a medio camino entre el código fuente y el código máquina. Especialmente en el caso del IR de LLVM utilizado en este trabajo, esta fase nos permite tener una mejor base para realizar diferentes optimizaciones y conversiones a diferentes formas de lenguaje máquina, según la arquitectura objetivo.
    \item [Optimización]: Esta fase consisten en realizar cambios o mejoras en la forma del lenguaje, sin alterar el significado original. Estas pueden ser eliminaciones de código no alcanzable, redefiniciones de algunas estructuras para evitar generar más variables de las necesarias, gestión eficiente de la memoria, etc.
    \item [Generación de código final]: Esta es la fase final del proceso de compilación, durante este paso se convierte el IR a código máquina, el cual está totalmente listo para ser ejecutado.
\end{description}
\capitulo{4}{Técnicas y herramientas}

Esta parte de la memoria tiene como objetivo presentar las técnicas metodológicas y las herramientas de desarrollo que se han utilizado para llevar a cabo el proyecto. Si se han estudiado diferentes alternativas de metodologías, herramientas, bibliotecas se puede hacer un resumen de los aspectos más destacados de cada alternativa, incluyendo comparativas entre las distintas opciones y una justificación de las elecciones realizadas. 
No se pretende que este apartado se convierta en un capítulo de un libro dedicado a cada una de las alternativas, sino comentar los aspectos más destacados de cada opción, con un repaso somero a los fundamentos esenciales y referencias bibliográficas para que el lector pueda ampliar su conocimiento sobre el tema.

\section{Metodología de desarrollo}
El desarrollo del proyecto se ha organizado siguiendo una metodología ágil, concretamente el marco de trabajo \textbf{SCRUM}, apoyado en la herramienta \textbf{ZenHub} para su gestión y seguimiento. La elección de una metodología ágil se justifica por su \textit{flexibilidad} y \textit{adaptabilidad}, características especialmente valiosas en un proyecto de investigación y desarrollo como este, donde la evolución del conocimiento y del propio software ocurre de forma simultánea e iterativa.

El flujo de trabajo se estructura mediante un tablero \textbf{Kanban} en ZenHub, el cual permite visualizar de manera clara el estado actual del proyecto. En dicho tablero se gestionan las \textit{issues} y los \textit{milestones} de GitHub, representando las distintas tareas, etapas o funcionalidades en desarrollo.  
Cada elemento del tablero incluye información relevante como la prioridad, la duración estimada, el tipo de tarea y su relación con otras, lo que facilita la planificación y el control del progreso de manera visual y dinámica.

Esta metodología fomenta un desarrollo incremental y una mejora continua del código, permitiendo incorporar nuevas funcionalidades, corregir errores y ajustar los objetivos a medida que avanza el trabajo, garantizando así un resultado más coherente y alineado con los objetivos técnicos del proyecto.

\section{Metodología de programación}
El desarrollo del compilador y de las herramientas asociadas se ha abordado desde el paradigma de la \textbf{programación orientada a objetos} (OOP, por sus siglas en inglés). Este enfoque resulta especialmente adecuado para un proyecto de estas características, ya que facilita la modularización del código y la definición clara de las responsabilidades de cada componente del sistema (analizador léxico, analizador sintáctico, generador de código, etc.).

La OOP favorece la reutilización, extensibilidad y mantenibilidad del software, permitiendo aislar los distintos módulos del compilador y facilitar su evolución futura. Además, este paradigma se integra de forma natural con las herramientas empleadas, como \textbf{ANTLR} y \textbf{LLVM}, ambas diseñadas con arquitecturas orientadas a objetos.

Otro motivo determinante en la elección de este enfoque es el uso de \textbf{C++} como lenguaje de implementación. C++ ofrece un modelo de programación fuertemente orientado a objetos, combinado con una gestión eficiente de recursos y un rendimiento elevado, aspectos cruciales en el desarrollo de compiladores y sistemas de bajo nivel. La compatibilidad nativa de C++ con las bibliotecas de ANTLR y LLVM refuerza aún más la idoneidad de esta metodología para el proyecto.


\section{Elección de herramientas}
A continuación se presentan las principales herramientas empleadas en el desarrollo del proyecto, así como la justificación de su elección frente a otras alternativas disponibles. Las herramientas seleccionadas se han escogido teniendo en cuenta su madurez, soporte, documentación, integración con C++ y adecuación a los objetivos técnicos del proyecto.

\subsection{ANTLR}
ANTLR (Another Tool for Language Recognition) es una herramienta ampliamente utilizada para la generación automática de analizadores léxicos y sintácticos a partir de gramáticas formales. Su principal fortaleza radica en la facilidad con la que permite definir y mantener gramáticas complejas mediante una sintaxis clara y expresiva, generando código eficiente y legible para distintos lenguajes de programación, incluido C++.

Entre sus características más destacables se encuentran:
\begin{itemize}
    \item Soporte nativo para gramáticas LL(*) que permiten manejar ambigüedades complejas.
    \item Generación automática de analizadores léxicos y sintácticos a partir de un único archivo de gramática.
    \item Aunque está diseñado para usarser con Java, actualmente cuenta con compatibilidad para muchos lenguajes de destino, entre ellos C++.
\end{itemize}

\textbf{Alternativas consideradas:}
\begin{itemize}
    \item \textbf{Bison y Flex:} Son las herramientas clásicas en la generación de analizadores. Sin embargo, en el contexto elegido ANTLR cuenta con una estructura orientada a objetos más integrable con otras herramientas.
    \item \textbf{PEG (Parsing Expression Grammar):} Ofrecen gran expresividad y flexibilidad, pero carecen del soporte multiplataforma y de la madurez de ANTLR, además de no integrarse de forma fácil con C++.
    \item \textbf{JavaCC:} Alternativas orientadas a Java, no adecuadas para este proyecto al estar centrado en C++.
\end{itemize}

\textbf{Justificación de la elección:}  
ANTLR se considera la opción más equilibrada entre facilidad de uso, potencia y soporte técnico. Permite definir la gramática del lenguaje de manera declarativa y coherente, integrándose perfectamente con el entorno C++. Además, su ecosistema y documentación la convierten en una herramienta moderna y sólida frente a las opciones más tradicionales.

\subsection{LLVM}
LLVM (Low Level Virtual Machine) es una infraestructura modular y extensible para el desarrollo de compiladores, enlazadores y optimizadores de código. Su arquitectura se basa en una representación intermedia (IR) que actúa como un lenguaje ensamblador universal, lo que permite desarrollar compiladores portables, eficientes y escalables.

Entre sus principales ventajas destacan:
\begin{itemize}
    \item \textbf{Representación intermedia (IR):} Permite una separación clara entre la fase de análisis y la de generación de código, haciendo posible optimizaciones y traducciones a múltiples arquitecturas.
    \item \textbf{Portabilidad:} LLVM IR puede compilarse a diferentes arquitecturas (x86, ARM, RISC-V, etc.) sin modificar el compilador fuente.
    \item \textbf{Integración con C++:} Al estar implementado en C++, se integra de manera natural con el entorno de desarrollo del proyecto.
    \item \textbf{Extensibilidad:} Su diseño modular permite hacer cambios como optimizaciones, cambio de backend o cambios en herramientas de análisis de forma sencilla.
\end{itemize}

\textbf{Alternativas consideradas:}
\begin{itemize}
    \item \textbf{GCC (GNU Compiler Collection):} Ofrece una infraestructura madura, pero su API interna no está diseñada para un uso externo, lo que dificulta su integración en proyectos ajenos al entorno GNU.
    \item \textbf{QBE Compiler Infrastructure:} Es más ligera que LLVM, pero carece del ecosistema, documentación y soporte de comunidad que caracterizan a LLVM.
    \item \textbf{Cranelift:} Un backend moderno y rápido, pero enfocado en compilación JIT (just-in-time) y no tan versátil para la generación de IR o compiladores tradicionales.
\end{itemize}

\textbf{Justificación de la elección:}  
LLVM ofrece el equilibrio ideal entre potencial, documentación y flexibilidad, su diseño modular permite construir un compilador completamente funcional sin preocuparse por la dependencia de la arquitectura de destino. Además, su integración con C++ simplifica enormemente la generación de IR y la gestión de optimizaciones.
\capitulo{5}{Aspectos relevantes del desarrollo del proyecto}

El comienzo del proyecto parte del objetivo de conseguir una versión MVP (Minimal Viable Product) que pudiese aportar una base sólida y fácil de ampliar. En este caso el MVP se enfocó en la creación de un micro compilador capaz de procesar expresiones aritméticas y lógicas, puesto que esta estructura es básica y está presente en todos los lenguajes de programación.

La primera fase de este desarrollo comienza con la creación de un lexer donde se definen algunas estructuras básicas como los lierales de número entero, número con decimales o las cadenas de caracteres, así como los operadores que se van a emplear para estas expresiones. Una vez obtenido el lexer el siguiente paso es construir un parser capaz de reconocer estructuras gramaticales más complejas donde se conectan el uso de los operadores y los literales definidos anteriormente.

A continuación se debe implementar una estructura de datos que sea capaz de representar cualquier tipo de estructura sintáctica bien formada de este lenguaje, la estructura del árbol de sintaxis abstracta (AST por sus siglas en inglés) será la encargada de crear este representacion en tiempo de ejecución del compilador, almacenando los datos relevantes de cada estructura para su posterior análisis en las siguientes fases.
Para formar este AST se empleará un visitor patterns (patrón de visita), objeto que recibe el nombre de ASTBuilder, el cual se encarga de visitar cada contexto generado de forma automática por el parser de ANTLR, dando por resultado unos nodos enlazados, el AST. Adicionalmente, ASTBuilder será el encargado de generar una representación visual del árbol a medida que lo va construyendo, para conseguir esto se escriben cada nodo en el formato de la biblioteca `forest` de LaTex a medida que estos se van formando, garantizando su orden y asociación correctos.



\capitulo{6}{Trabajos relacionados}

Este apartado sería parecido a un estado del arte de una tesis o tesina. En un trabajo final grado no parece obligada su presencia, aunque se puede dejar a juicio del tutor el incluir un pequeño resumen comentado de los trabajos y proyectos ya realizados en el campo del proyecto en curso. 
\capitulo{7}{Conclusiones y Líneas de trabajo futuras}

Todo proyecto debe incluir las conclusiones que se derivan de su desarrollo. Éstas pueden ser de diferente índole, dependiendo de la tipología del proyecto, pero normalmente van a estar presentes un conjunto de conclusiones relacionadas con los resultados del proyecto y un conjunto de conclusiones técnicas. 
Además, resulta muy útil realizar un informe crítico indicando cómo se puede mejorar el proyecto, o cómo se puede continuar trabajando en la línea del proyecto realizado. 

% Quitar luego, evita errores de compilacion
\cite{koza92}

\bibliographystyle{plain}
\bibliography{bibliografia}

\end{document}